\subsection{Interpretation}
  \label{subsec:continuum-interpretation}

  The continuum scaling properties confirm that
  the Newtonian kernel arises from the elliptic nature
  of the operator rather than from lattice-specific effects.
  The discrete simulations reproduce the expected continuum behavior,
  while making explicit the role of finite spatial support
  in generating a macroscopic amplitude shift.

  Within this framework, gravity appears as a long-range
  resolvent-mediated response of a projective entropy functional.
  The derivation is restricted to the weak-field,
  Newtonian regime and does not rely on additional geometric structure.

  The numerical threshold observed in Section~4.3 admits
  a natural spectral interpretation.
  Since the entropy variation is governed by
  $\delta S_\Pi = \frac12 \mathrm{Tr}(A^{-1}\delta A)$,
  the long-range response is controlled by the infrared-enhanced
  weight $1/\lambda$ in the spectral sum.
  A perturbation activates the macroscopic $1/r$ channel
  only if it overlaps sufficiently with low-lying modes.
  In scaling form, this may be expressed as
  \begin{equation}
    \int_{V_R} \delta A(x)\,dx \gtrsim \mathcal{C}_{\mathrm{IR}},
  \end{equation}
  where $\mathcal{C}_{\mathrm{IR}}$ denotes an infrared stiffness scale
  determined by the low-frequency spectral density of~$A$.
  Below this threshold, the entropy variation is non-zero
  but remains spectrally local.

  This supports an operational reinterpretation of mass.
  Gravitational mass is not an intrinsic point-like charge,
  but the asymptotic amplitude of the resolvent-mediated response.
  Only perturbations whose spatial support exceeds
  a coherence radius $R_c \sim \Lambda_{\mathrm{IR}}^{-1/2}$
  produce a stable macroscopic $1/r$ profile.
  In this sense, mass measures spectral coherence
  rather than local energy density alone.

  The activation threshold also acts as an infrared filtering mechanism.
  Idealized point-like sources fail to excite the non-local channel
  because they do not satisfy the coherence condition.
  The singular limit of a point perturbation therefore does not
  automatically translate into a divergent macroscopic response.
  The long-range interaction is restricted to perturbations
  carrying sufficient infrared spectral weight.

  These results suggest that Newtonian gravity can be understood
  as a collective spectral phenomenon:
  the $1/r$ law follows from ellipticity in three dimensions,
  while its observable amplitude is conditioned
  by infrared coherence of the source.
